\documentclass[11pt]{article}
\usepackage[margin=1in]{geometry}
\usepackage{amsmath,amssymb,amsthm}
\usepackage{hyperref}
\newtheorem{theorem}{Theorem}
\newtheorem{lemma}[theorem]{Lemma}
\newtheorem{proposition}[theorem]{Proposition}
\newtheorem{corollary}[theorem]{Corollary}
\theoremstyle{definition}
\newtheorem{definition}[theorem]{Definition}
\theoremstyle{remark}
\newtheorem{remark}[theorem]{Remark}

\title{Disproof of Conjecture 00000007145:\\
a 3-dimensional polytope with a non-Hamiltonian graph}
\author{jilint777}
\date{2026-10-04}

\begin{document}
\sloppy
\maketitle

\begin{abstract}
Conjecture 00000007145 asserts that non-Hamiltonian graphs of 4-dimensional
polytopes exist and that \emph{the minimal dimension of such examples is four}.
The minimality clause is false. The rhombic dodecahedron
$\operatorname{conv}\{(\pm2,0,0),(0,\pm2,0),(0,0,\pm2),(\pm1,\pm1,\pm1)\}$ is a
3-dimensional polytope. Its graph is bipartite with parts of sizes $8$ and $6$,
so it has no Hamiltonian cycle (and not even a Hamiltonian path). Because every
polygon has a Hamiltonian graph, the true minimal dimension is $3$. The first
clause is true: the pyramid over the rhombic dodecahedron is a 4-polytope with a
non-Hamiltonian graph. Polytopes, vertices, edges, dimension and Hamiltonian
cycles are defined from scratch in Lean~4 (core library only). Lean proves that the
rhombic dodecahedron is a 3-polytope, determines its whole graph, proves
that the graph is not Hamiltonian, and proves $\neg$\,\texttt{MinDimIsFour}.
\end{abstract}

\section{The conjecture and its reading}
\begin{quote}
\emph{Definition:} The Hamiltonicity of polytope graphs with skeleton cycles.
\emph{Conjecture:} Non-Hamiltonian examples of four-dimensional polytope graphs
exist, and the minimal dimension of the examples is four. (minimal
four-dimensional non-Hamiltonian skeletons)
\end{quote}
The Chinese version says the same: there are non-Hamiltonian examples of graphs of
4-dimensional polytopes, and the minimal dimension of the examples is four.

The \emph{graph} (1-skeleton) of a convex polytope $P$ has the vertices of $P$ as
its vertices and the edges (1-dimensional faces) of $P$ as its edges. It is
\emph{Hamiltonian} if it has a cycle through every vertex exactly once. An
\emph{example} is a polytope whose graph is not Hamiltonian. The conjecture is the
conjunction of
\begin{itemize}
\item[(A)] there is a 4-dimensional polytope with a non-Hamiltonian graph;
\item[(B)] the minimal dimension of a polytope with a non-Hamiltonian graph is $4$;
that is, every polytope of dimension $<4$ has a Hamiltonian graph.
\end{itemize}
We refute (B). This refutes the conjunction. We also prove (A), so the conjecture
fails exactly through its minimality clause.

\paragraph{Readings.}
\begin{enumerate}
\item \emph{Degenerate polytopes.} A point (dimension $0$) and a segment
(dimension $1$) have graphs $K_1$ and $K_2$, which contain no cycle at all. Under
the usual convention (a Hamiltonian cycle needs at least $3$ vertices) these are
non-Hamiltonian, so the minimal dimension would trivially be $0$. We do not use this.
We read (B) charitably as ``every polytope of dimension $2$ or $3$ has a
Hamiltonian graph''. Our counterexample has dimension exactly $3$, so it refutes
(B) under every convention for dimensions $0,1$. It even refutes the weakest
consequence of (B): ``every 3-polytope has a Hamiltonian graph''.
\item \emph{Cycle or path.} The graph of our example has no Hamiltonian cycle
and also no Hamiltonian path (Remark~\ref{rem:path}). So (B) also fails if
``Hamiltonian'' meant ``traceable''.
\item \emph{Abstract graphs.} If ``examples'' are abstract non-Hamiltonian graphs
and ``dimension'' is the least dimension of a polytope realising the graph, (B)
still fails: our graph is the graph of a 3-polytope, so that least dimension is at
most $3$.
\item \emph{Simple polytopes.} If only simple polytopes (every vertex on exactly $d$
facets) were meant, (B) still fails. Tutte's graph \cite{tutte} is cubic, planar
and 3-connected. By Steinitz's theorem it is the graph of a simple 3-polytope, and
it is not Hamiltonian. This reading is not formalized; we only cite it.
\item \emph{Simplicial polytopes.} If only simplicial polytopes (all facets are
simplices) were meant, (B) still fails in dimension $3$. The triakis octahedron
\[T=\operatorname{conv}\{(\pm12,0,0),(0,\pm12,0),(0,0,\pm12),(\pm5,\pm5,\pm5)\}\]
has
$14$ vertices, $36$ edges and $24$ triangular facets. Its $8$ apexes
$(\pm5,\pm5,\pm5)$ are pairwise non-adjacent and $8>14/2$, so by
Lemma~\ref{lem:indep} its graph is not Hamiltonian. The script
\texttt{verify.py} checks the facets, the edges, that every facet is a triangle, and
non-Hamiltonicity by exhaustive search. This example is not formalized in Lean.
\item \emph{Facet adjacency (dual graph).} If ``polytope graph'' meant the graph
whose nodes are the facets, two facets being adjacent when they share a ridge, (B)
still fails in dimension $3$. The cuboctahedron, the convex hull of all
permutations of $(\pm1,\pm1,0)$, is the polar dual of
the rhombic dodecahedron $R$ of Section~4. Its $8$ triangles and $6$ squares are its
facets, and every edge lies on one triangle and one square. Its facet-adjacency
graph is therefore the graph of $R$, which is not Hamiltonian
(Theorem~\ref{thm:main}). \texttt{verify.py} computes this dual graph directly and
checks that it is not Hamiltonian.
\item \emph{Tautological reading.} One could read ``the examples'' as only the
4-dimensional examples of the first clause. Then ``their minimal dimension is
four'' is true by definition and says nothing beyond (A), and the minimality clause
would be empty. So this cannot be the intended meaning. The meaningful reading, used
throughout, takes the minimum over all polytopes with non-Hamiltonian graphs.
\end{enumerate}

\section{Definitions}\label{sec:defs}
Let $V\subset\mathbb{R}^d$ be finite and $P=\operatorname{conv}V$. A
\emph{face} of $P$ is a set $\{x\in P: c\cdot x=\max_P c\}$ for some
$c\in\mathbb{R}^d$. The maximum of $c$ over $P$ equals its maximum over $V$, and the
face is the convex hull of the maximising points of $V$. We always take $V$ to be
the vertex set of $P$, so that every point of $V$ is a vertex. Then:
\begin{itemize}
\item $p\in V$ is a \emph{vertex} iff some $c$ satisfies $c\cdot q<c\cdot p$ for
all $q\in V\setminus\{p\}$;
\item $\{p,q\}\subseteq V$, $p\neq q$, is an \emph{edge} iff some $c$ satisfies
$c\cdot p=c\cdot q$ and $c\cdot r<c\cdot p$ for all $r\in V\setminus\{p,q\}$
(the face is then exactly the segment $[p,q]$);
\item $\dim P=d$ iff some $d+1$ points $p,b_1,\dots,b_d\in V$ are affinely
independent, i.e. $\det(b_1-p,\dots,b_d-p)\neq0$. (Every $d$-polytope is
affinely isomorphic to a full-dimensional polytope in $\mathbb{R}^d$, and affine
isomorphisms preserve the face lattice, hence the graph.)
\end{itemize}

\begin{lemma}[integer functionals suffice]\label{lem:int}
Let $V\subset\mathbb{Q}^d$ be finite. In the definitions of vertex and edge above,
the functional $c$ may be taken in $\mathbb{Z}^d$.
\end{lemma}
\begin{proof}
For an edge $\{p,q\}$ let $L=\{x\in\mathbb{R}^d: x\cdot(p-q)=0\}$; for a vertex let
$L=\mathbb{R}^d$. $L$ is the kernel of a rational matrix, so it has a basis
$u_1,\dots,u_m$ of rational vectors. Write $c=\sum\lambda_iu_i$. The finitely many
strict inequalities $(\sum\mu_iu_i)\cdot(p-r)>0$ define an open subset of
$\mathbb{R}^m$ that contains $\lambda$. So it contains a rational point $\mu$, and
$c'=\sum\mu_iu_i\in L\cap\mathbb{Q}^d$ satisfies the same conditions as $c$.
Multiplying $c'$ by a positive common denominator gives an integer functional.
\end{proof}
So for polytopes with rational vertices the integer definitions used in Lean
describe exactly the usual (real) vertices and edges. A rational polytope can be
scaled to have integer vertices without changing its graph.

\begin{definition}
A \emph{Hamiltonian cycle} of the graph of $P$ is a sequence
$v_0,v_1,\dots,v_{n-1}$ listing every vertex exactly once, with $n\ge3$, such that
$v_i v_{i+1}$ (indices mod $n$) are edges.
\end{definition}

\section{Certificates for edges and non-edges}
\begin{lemma}[non-edge certificate]\label{lem:cert}
Let $p\ne q$ in $V$. Suppose there are an integer $k>0$, points
$w_1,\dots,w_s\in V\setminus\{p,q\}$ and integers $m_i\ge0$ with
$\sum m_i=2k$ and $k(p+q)=\sum m_iw_i$. Then $\{p,q\}$ is not an edge.
\end{lemma}
\begin{proof}
Suppose $c$ satisfies $c\cdot p=c\cdot q=t$ and $c\cdot r<t$ for all other
$r\in V$. Then
\[2kt=c\cdot k(p+q)=\sum m_i\,c\cdot w_i<\sum m_i\,t=2kt,\]
where the inequality is strict because $\sum m_i=2k>0$. This is a contradiction. The
argument uses only the arithmetic of an ordered ring, so it applies to real $c$
directly (Lemma~\ref{lem:int} is not even needed here). Geometrically, the
midpoint of $[p,q]$ is the convex combination $\sum\frac{m_i}{2k}w_i$ of the
other vertices.
\end{proof}

\begin{lemma}[independence criterion]\label{lem:indep}
Let $G$ be a graph on $n$ vertices and let $I$ be a set of pairwise non-adjacent
vertices with $|I|>n/2$. Then $G$ has no Hamiltonian cycle.
\end{lemma}
\begin{proof}
In a Hamiltonian cycle $v_0,\dots,v_{n-1}$ the successor map
$v_i\mapsto v_{i+1}$ is injective. It sends every vertex of $I$ to a vertex
outside $I$, because $I$ is independent. Hence $|I|\le n-|I|$, a contradiction.
\end{proof}
The Lean proof of Lemma~\ref{lem:indep} is a list induction. If no two consecutive
entries of a list are both in $I$, then the number of entries in $I$ is at most
the number outside $I$, plus one if the list starts in $I$. A cycle that starts in
$I$ is rotated by one step first.

\section{The rhombic dodecahedron}
Let
\[V=\underbrace{\{(\pm2,0,0),(0,\pm2,0),(0,0,\pm2)\}}_{O\ (6\text{ axis points})}
\cup\underbrace{\{(\pm1,\pm1,\pm1)\}}_{C\ (8\text{ cube points})},\qquad
R=\operatorname{conv}V.\]
This is the rhombic dodecahedron. Its $12$ facets are the rhombi
$\pm x_i\pm x_j=2$ ($i<j$), for example
$\{(2,0,0),(1,1,1),(0,2,0),(1,1,-1)\}\subset\{x+y=2\}$.

\begin{proposition}\label{prop:R}
\begin{enumerate}
\item Every point of $V$ is a vertex of $R$, and $\dim R=3$.
\item The edges of $R$ are exactly the $24$ pairs
$\{s,\,2s_ie_i\}$ with $s\in C$ and $i\in\{1,2,3\}$, where $e_i$ is the $i$-th unit vector. Every edge joins
a point of $C$ to a point of $O$.
\end{enumerate}
\end{proposition}
\begin{proof}
(1) For $p\in V$ take $c=p$. Then $c\cdot p=|p|^2\in\{3,4\}$. For $q\ne p$ in $V$ we
have $p\cdot q\le 2<3$ if $p\in C$, and $p\cdot q\le 2<4$ if $p\in O$. The points
$(0,0,-2),(2,0,0),(0,2,0),(0,0,2)$ give
$\det\bigl((2,0,2),(0,2,2),(0,0,4)\bigr)=16\ne0$.

(2) \emph{Edges.} For $s\in C$ and $a=2s_ie_i$ take $c=s+s_ie_i$. Then
$c\cdot s=c\cdot a=4$, and every other point $r\in V$ has $c\cdot r\le2$: for the
axis points $\pm2e_j$ with $j\ne i$ the value is $\pm2$, for $-2s_ie_i$ it is $-4$,
and for a cube point $t\ne s$ it is $s\cdot t+s_it_i\le1+1=2$.

\emph{Non-edges.} The remaining $\binom{14}{2}-24=67$ pairs have certificates as
in Lemma~\ref{lem:cert}, all with $k\in\{1,2\}$, for example:
\begin{align*}
(1,1,1)+(1,1,-1)&=(2,0,0)+(0,2,0) &&\text{(cube points differing in one sign)},\\
2\bigl((1,1,1)+(1,-1,-1)\bigr)&=2\cdot(2,0,0)+(0,2,0)+(0,-2,0) &&\text{(differing in two signs)},\\
(1,1,1)+(-1,-1,-1)&=(2,0,0)+(-2,0,0) &&\text{(opposite)},\\
(2,0,0)+(0,2,0)&=(1,1,1)+(1,1,-1) &&\text{(two axis points)},\\
(2,0,0)+(-2,0,0)&=(0,2,0)+(0,-2,0) &&\text{(opposite axis points)},\\
(2,0,0)+(-1,1,1)&=(0,2,0)+(1,-1,1) &&\text{(non-incident axis and cube points)}.
\end{align*}
All $67$ certificates are listed in \texttt{lean4/Main.lean} (\texttt{rdNonEdges}),
checked there by \texttt{decide}, and re-checked by \texttt{verify.py}. Up to
the symmetries of the cube, every non-edge is of one of the six types above.
\end{proof}

\begin{theorem}\label{thm:main}
The graph of the 3-polytope $R$ has no Hamiltonian cycle. Hence clause (B) is
false, and Conjecture 00000007145 is false.
\end{theorem}
\begin{proof}
By Proposition~\ref{prop:R}, the $8$ cube points are pairwise non-adjacent, and
$8>14/2$. Apply Lemma~\ref{lem:indep}.
\end{proof}

\begin{remark}\label{rem:path}
The graph of $R$ is bipartite with parts $C$ and $O$ of sizes $8$ and $6$. A path
alternates between the parts, so the numbers of its vertices in the two parts
differ by at most $1$. Hence there is not even a Hamiltonian path.
\end{remark}

\begin{corollary}
The minimal dimension $d\ge2$ of a polytope with a non-Hamiltonian graph is $3$.
\end{corollary}
\begin{proof}
A convex $n$-gon ($n\ge3$) has the $n$-cycle as its graph, which is Hamiltonian.
Theorem~\ref{thm:main} gives a 3-dimensional example.
\end{proof}

\begin{remark}
Other classical 3-dimensional examples are the triakis octahedron (for example,
$(\pm12,0,0),(0,\pm12,0),(0,0,\pm12)$ together with $(\pm5,\pm5,\pm5)$: $14$ vertices,
of which the $8$ apexes are independent), the Herschel enneahedron ($11$ vertices,
the smallest non-Hamiltonian polyhedral graph \cite{herschel}) and Tutte's simple
polyhedron \cite{tutte}. \texttt{verify.py} also checks the triakis octahedron.
\end{remark}

\section{The first clause: a 4-dimensional example}
Let $\Pi=\operatorname{conv}\bigl(\{(p,0):p\in V\}\cup\{(0,0,0,1)\}\bigr)\subset\mathbb{R}^4$
be the pyramid over $R$.
Every listed point is a vertex (again $c=p$ works: $|p|^2\in\{1,3,4\}$, while
$p\cdot q\le2$ for $q\neq p$ and the apex pairs to $0$ with base points). The points
$(0,0,-2,0)$, $(2,0,0,0)$, $(0,2,0,0)$, $(0,0,2,0)$ and $(0,0,0,1)$ are affinely independent,
so $\dim\Pi=4$. For two distinct cube points of the base, the lifted certificate of
Proposition~\ref{prop:R} (append a $0$ coordinate) satisfies Lemma~\ref{lem:cert}
in $\Pi$. So the $8$ lifted cube points are pairwise non-adjacent in the graph
of $\Pi$, and $8>15/2$. By Lemma~\ref{lem:indep}, $\Pi$ is a 4-polytope with a
non-Hamiltonian graph, and (A) holds. (The facet computation in \texttt{verify.py}
gives $13$ facets and $38$ edges.)

\section{Lean formalization}
\texttt{lean4/Main.lean} is a self-contained Lean~4.19.0 project (core library
only, no Mathlib). All objects are defined from scratch:
\begin{itemize}
\item Points and functionals are integer vectors \texttt{Pt := List Int}, with
\texttt{dot}, \texttt{vadd}, \texttt{smul} and \texttt{lincomb}, together with the
linearity lemmas \texttt{dot\_vadd}, \texttt{dot\_smul} and \texttt{dot\_lincomb}.
\texttt{det} is the Laplace-expansion determinant, with sanity examples.
\item \texttt{IsVertex d V p}, \texttt{Adj d V p q} (edge), \texttt{FullDim d V}
are the definitions of Section~\ref{sec:defs} with $c\in\mathbb{Z}^d$
(justified by Lemma~\ref{lem:int}). \texttt{structure Polytope} bundles a dimension
$d$, a vertex list in $\mathbb{Z}^d$ without repetitions, a proof that every listed
point is a vertex, and a proof of full dimensionality.
\item \texttt{HamCycle d V cyc}: \texttt{cyc} has no repetitions, has the same
elements as \texttt{V}, has length $\ge3$, and \texttt{Path (Adj d V) (cyc ++
cyc.take 1)}, i.e.\ cyclically consecutive entries are adjacent.
\texttt{Polytope.Hamiltonian P} means $\exists$ \texttt{cyc}, \texttt{HamCycle}.
\item \texttt{MinDimIsFour} is (A) $\wedge$ ``every polytope with $2\le d$ and a
non-Hamiltonian graph has $d\ge4$''.
\end{itemize}
Main theorems:
\begin{itemize}
\item \texttt{not\_adj\_of\_cert} (Lemma~\ref{lem:cert}) and \texttt{no\_hamCycle}
(Lemma~\ref{lem:indep}) are general, for any $d$ and $V$.
\item The rhombic dodecahedron is \texttt{RD : Polytope}, with \texttt{RD.d = 3}.
The theorem \texttt{rd\_adj\_iff} determines the full graph of $R$: the 24 edges come with exposing functionals
(\texttt{rd\_edges\_ok}), the 67 non-edges come with certificates
(\texttt{rd\_nonedges\_ok}), and \texttt{rd\_cover} shows that the two tables
cover every pair. \texttt{rd\_edges\_bipartite} and \texttt{cube\_independent} follow.
\item \texttt{RD\_not\_hamiltonian : $\neg$ RD.Hamiltonian};
\texttt{exists\_nonHamiltonian\_3polytope}.
\item The main theorem is \texttt{conjecture\_00000007145\_false :}
\texttt{$\neg$ MinDimIsFour}. The variant \texttt{conjecture\_00000007145\_false'}
covers the reading without the restriction $d\ge2$.
\item \texttt{exists\_nonHamiltonian\_4polytope} (the pyramid \texttt{Pyr}) proves (A).
\item Non-vacuity: \texttt{Oct\_hamiltonian} shows that the octahedron is a
\texttt{Polytope} with a Hamiltonian cycle, and \texttt{oct\_nonedge} shows that
\texttt{Adj} can fail. Together with \texttt{rd\_adj\_iff}, this shows the
predicates are neither empty nor trivial.
\end{itemize}
There is no \texttt{sorry}, no \texttt{native\_decide}, and no added axioms.
\texttt{\#print axioms} reports only \texttt{propext}, \texttt{Classical.choice}
and \texttt{Quot.sound}.

\paragraph{Limitations (stated precisely).} Lean polytopes have integer
vertices and integer functionals. This is enough for the disproof: integer
polytopes are polytopes, so refuting ``all integer 3-polytopes are Hamiltonian''
refutes ``all 3-polytopes are Hamiltonian''. By Lemma~\ref{lem:int} (proved here,
not in Lean), the integer notions of vertex and edge agree with the real ones.
The disproof itself needs only the certificate lemma (\texttt{not\_adj\_of\_cert},
Lemma~\ref{lem:cert}). Its proof holds verbatim for real functionals, so the Lean
non-adjacency facts, and hence non-Hamiltonicity, carry over to real polytopes
directly. Lemma~\ref{lem:int} is needed only to identify the $24$ positive edges
of $R$ as real edges. The
Hamiltonian-path reading (Remark~\ref{rem:path}) and the simple-polytope reading
(Tutte) are argued in this report only; \texttt{verify.py} also checks the path
statement by exhaustive search.

\section{Independent check}
\texttt{verify.py} (Python~3 standard library) uses a different method. It
finds all facets by brute force over $d$-subsets with exact integer hyperplanes.
A point is a vertex iff the intersection of the facets containing it is the point
itself. A pair $\{p,q\}$ is an edge iff the intersection of the facets containing
both is exactly $\{p,q\}$. For $R$ it confirms $14$ vertices, $12$ rhombic facets,
$24$ edges ($V-E+F=2$) and bipartiteness, and an exhaustive depth-first search
finds no Hamiltonian cycle and no Hamiltonian path. It re-checks every edge
functional and every non-edge certificate of the Lean tables against the
facet-computed graph. For the pyramid it confirms dimension $4$, $13$ facets,
$38$ edges and no Hamiltonian cycle. It finds a Hamiltonian cycle in the octahedron
(a positive control). It checks that the triakis octahedron is simplicial and
non-Hamiltonian, and that the facet-adjacency graph of the cuboctahedron is not
Hamiltonian.

\section{Reproduction}
\begin{verbatim}
cd lean4 && lake build     # Lean 4.19.0, core only
cd .. && python3 verify.py # Python 3 standard library
pdflatex report.tex && pdflatex report.tex
\end{verbatim}

\begin{thebibliography}{9}
\bibitem{gr} B.~Gr\"unbaum, \emph{Convex Polytopes}, 2nd ed., Springer, 2003.
\bibitem{zi} G.~M.~Ziegler, \emph{Lectures on Polytopes}, Springer, 1995.
\bibitem{tutte} W.~T.~Tutte, On Hamiltonian circuits, \emph{J. London Math. Soc.}
21 (1946), 98--101.
\bibitem{herschel} A.~S.~Herschel, Sir Wm.\ Hamilton's Icosian Game,
\emph{Quart. J. Pure Appl. Math.} 17 (1880), 305.
\end{thebibliography}
\end{document}
